\chapter{Aggiunzioni tra categorie}\label{cap_aggiunti}
\epigraph{L'organizzazione hegeliana della conoscenza si realizza concretamente in termini di funtori aggiunti su di una qualsiasi categoria matematica, e viene utilizzata
	per dare una definizione precisa della dimensione necessaria per una rappresentazione.}{F.\ W.\ Lawvere, \cite{Lawvere1989Taco}}
\subsubsection*{Esercizi}
\begin{enumerate}
	\item
	\item
	\item
	\item
	\item
\end{enumerate}

\section{Addenda \& paralipomena}\color{gray}
\paragraph{Le aggiunzioni relative}\label{sec_agg_rel}

\begin{color}{blue}
La nozione di aggiunzione relativa non \`e certamente cos\`i centrale come quella di aggiunzione e probabilmente 
non sar\`a sviluppata nel libro. Qua propongo solo di darne la definizione e l'ovvia propriet\`a di preservazione dei 
colimiti perch\'e \`e particolarmente utile per semplificare la presentazione del fatto che una categoria algebrica
\(\ctAlgT\) \`e il completamento libero della teoria algebrica \(\ctT\) rispetto ai colimiti setacciati (proposizione 
\ref{prop_univ_yon_alg}).
\end{color} 

\begin{definition}\label{def_agg_rel}
Un'aggiunzione relativa \`e data da un diagramma di categorie e funtori
\[ \xymatrix{ & \ctA \ar[d]^{i} \ar[ld]_{F} \\
\ctB \ar[r]_-{R} & \ctC } \]
e da una famiglia naturale di biiezioni
\( \{\Theta_{A,B} \colon \ctB(FA,B) \to \ctC(iA,RB)\}_{A \in \ctA, B \in \ctB} \).
\end{definition} 

\begin{remark}\label{oss_agg_rel_agg}
Un'aggiunzione \`e esattamente un'aggiunzione relativa in cui il funtore \(i \colon \ctA \fun \ctC\) \`e un'equivalenza.
\end{remark} 

La seguente proposizione generalizza il fatto che gli aggiunti sinistri preservano i colimiti. La dimostrazione \`e praticamente
identica a quella per le aggiunzioni e sar\`a omessa.

\begin{proposition}\label{prop_preserv_agg_rel}
Consideriamo un'aggiunzione relativa come nella definizione \ref{def_agg_rel} e un funtore \(H \colon \ctD \fun \ctA\).
Se \(\{\sigma_D \colon HD \to C\}_{D \in \ctD}\) \`e una famiglia naturale e se 
\[ \colim_{\ctD}(i \cmp H) = \langle iC, \{i(\sigma_D) \colon i(HD) \to iC\}_{D \in \ctD}\rangle \]
allora 
\[ \colim_{\ctD}(F \cmp H) = \langle FC, \{Fi(\sigma_D) \colon F(HD) \to FC\}_{D \in \ctD}\rangle \]
In particolare, \(F\) preserva tutti i colimiti che esistono in \(\ctA\) e che sono preservati da \(i\).
\end{proposition} 

\paragraph{Sui limiti e colimiti nelle sottocategorie riflessive}\label{sec_sottocat_rifl}

\begin{proposition}\label{prop_(co)lim_sottocat_rifl}
Consideriamo un'aggiunzione
$$\xymatrix{\ctA \ar@<-0.5ex>[r]_-{i} & \ctB \ar@<-0.5ex>[l]_{r}} \;\;\;\; r \dashv i$$
con unit\`a $\eta$ e counit\`a $\varepsilon.$ Assumiamo che $i$ sia pieno e fedele (per cui la counit\`a
$\varepsilon$ \`e un isomorfismo naturale) e sia $F \colon \ctD \fun \ctA$ un funtore definito su una categoria piccole $\ctD.$
\begin{enumerate}
\item Se $\colim(i \cdot F) = \langle C \in \ctB, \{s_D \colon i(FD) \to C\}_{D \in \ctD}\rangle,$ allora
$$\colim(F) = \langle r(C) \in \ctA,\{r(s_D) \cdot \varepsilon_{FD}^{-1} \colon FD \to r(i(FD)) \to r(C)\}_{D \in \ctD}\rangle$$
\item Se $\lim(i \cdot F) = \langle L \in \ctB, \{p_D \colon L \to i(FD)\}_{D \in \ctD}\rangle,$ allora
$$\lim(F) = \langle r(L) \in \ctA, \{q_D \colon r(L) \to FD\}_{D \in \ctD}\rangle$$
dove $q_D$ \`e la freccia che corrisponde \`a $p_D$ nella biiezione dell'aggiunzione $r \dashv i.$
\end{enumerate}
\end{proposition} 

\begin{remark}\label{oss_lim_sottocat_rifl}
Il punto 2 della propriet\`a \ref{prop_(co)lim_sottocat_rifl} merita un commento. Poich\'e il funtore
$i \colon \ctA \fun \ctB$ \`e un aggiunto destro, preserva i limiti. Quindi si ottiene
$$L = \lim(i \cdot F) \simeq i(\lim(F)) = i(rL)$$
Questo significa che il limite di $i \cdot F$ in realt\`a si trova gi\`a in $\ctA,$ ovvero che $\ctA$ \'e chiusa in $\ctB$
rispetto ai limiti.
\end{remark}

\paragraph{$\ctSet$ \`e cartesiana chiusa}\label{sec_set_cartcl}

\begin{color}{blue}
Questo \`e certamente uno degli esempi di base da mettere nel capitolo sui funtori aggiunti.
\end{color}

\begin{example}\label{ex_set_cartcl}
Per ogni insieme $Y,$ il funtore indotto
$$(-) \times Y \colon \ctSet \fun \ctSet$$
ha un aggiunto destro. Questa situazione si esprime dicendo che la categoria $\ctSet$ \`e cartesiana chiusa. 
L'aggiunto destro \`e il funtore rappresentabile $\ctSet(Y,-) \colon \ctSet \fun \ctSet$ e la counit\`a \`e data
dalla funzione di valutazione
$$\mathrm{ev} \colon \ctSet(Y,X) \times Y \to X \;\;\;\;\;\; \mathrm{ev}(f,y) = f(y)$$
\end{example}  

\color{black}